% Pursuing Stacks (cotes 134-2 to 134-8): archive cote 134-5, site page 6 (PDF page 7), typescript page 193.
% typescript. Draft compared: ps4.tex:249-307.
% AI-transcribed by Claude Opus 5.5 (claude-opus-5-5, adaptive thinking, effort high), Message Batches API, batch ps2 (msgbatch_01NVwr1Yuo8qJf5rykeGBXPK), 2026-11-12; not yet checked by eye.
\scanpage{134-5/6}
\mspage{193}

ment'' (``unfolding'') of stratified structures,  as a leading thread toward\unsure{s}
% NOTE: the right edge of the page is cut off; the last letters of several lines are hidden
a heuristic understanding of the natural stratification of \inserted{the} Mumford-Delign\unsure{e}
compac\unsure{ti}fications of modular multiplicities $M_{g,\hand{\nu}}$ (for curves of genus $g$
\unsure{andowed|endowed} with \hand{$\nu$} points).  This then led to the ``anabelian seminar'' which
\xout{\illegible{1}} turned short,  last spring.  Then a month or two sickness,  intensive
meditation \xout{\unsure{and}} for three or four months,  a few more months for settling
some important personal matter;  and now,since february,another mathemati%
-cal reflection started. \hand{$\rfloor$}I \xout{\unsure{have}} am unable to tell the meaning of this
alternation of periods of meditation on personal matters and periods of
mathematical reflection,  which has been going through my life for the
last seven years,  more and more,\xout{\unsure{like the}}- very much like the unceasing
up and down of waves,  or like a steady breathing going through my life,
without any attempt any longer of controlling it one way or other.  One
common moving force surely is the inborn curiosity  -  a thirst \over{of}{for} gettin\unsure{g}
acquainted with the juicy things of th\over{is}{e} \inserted[typed]{unexhaustible} world,  whether they be the brea-%
thing body of the \unsure{Beloved|beloved},  or the evasive substance of one's own life,
or the much less evasive substance of mathematical structures and their
delicate interplay.  This thirst in itself is of a nature quite different
of the ego's  -  it is the thirst of life to know about itself,  a primal
creative force which,one suspects,has been around forever,  long before
a human ego  -  a bundle \inserted{mainly} of fears,  of inhibitions and self-deceptions -
came into being.  Still I am aware that the ego is strongly involved in
the particular way in which the creative force expresses itself,  in my
own life or anyone's else's (when this force is allowed to come into pla\unsure{y}
at all ...).  The motivations behind any strong energy investment,  and
more particularly so when in it \over{in}{is} an activity attached with any kind of
social status or prestige, are a lot more complex and ego-\deleted{\unsure{loaded}}\inserted{driven} than
one generally cares to admit.  True,  ambition by itself is \deleted{\unsure{utterly}} power-%
less for discovering or understanding \inserted[typed]{or perceiving} \over{an}{an}ything substantial whatever,
neither a mathematical relationship nor the perfume of a flower.  In \deleted{the}
moments of work and \inserted{of} discovery,  in any creative moments in life, \inserted[typed]{as a matter of fact} ambitio\unsure{n}
is absent; the Artisan is a keen interest,  which is just one of the mani-%
fold aspects of love.  What \inserted[typed]{however} pushes us so relentlessly to work again,and
so often causes our life \inserted[typed]{and passion} gradually to dry off and become unsensitive,
even to the kindred passion of a fellow-being\hand{,} or to the un\unsure{sus}pected,
bea\inserted{u}ties and mysteries of the very field we are supposed to be ploughing
-  this force is neither love nor keen interest for the things \deleted{\unsure{an}}\inserted{and} beings
in this world.  It is interesting enough though,  and surely deserves
a close look !

\margin{6}
